% SGA 3，第 XI 篇；根据英文重建源翻译。
% 范围：仅注 1.6，止于注 1.7 之前。
% 权威来源：Polo--Gille Expo11.pdf，组件局部页 4--5；合卷页 726--727；印刷／原文页边页 132--133 和 88--89。
% Jacob Reinhold 的 CC BY 4.0 Markdown 谱系仅用于措辞比较；文字、公式、编号、注释及图均以 Polo--Gille PDF 为准。
% 精确续接／终点：注 1.7 之前。

\medskip\noindent\SGARefTarget{sga3:xi:remark:1:6}\textbf{注 1.6。} --- 当然，给定预概形态射 \(f:X\to Y\)，若 \(f\) 使 \(X\) 成为平滑（分别为非分歧、平展）\(Y\)-预概形，就称 \(f\) 为平滑（分别为非分歧、平展）态射。当 \(X\) 和 \(Y\) 为 \(S\)-预概形，且 \(f\) 为有限呈示 \(S\)-态射时，\(f\) 的这些性质可以直接用 \(f\) 定义的函子态射 \[
  F,G:(\mathbf{Sch})^\circ/S\longrightarrow(\mathbf{Ens})
\] 表达。也就是说，\(f\) 平滑（分别为非分歧、平展），若对每个在绝对意义下仿射的 \(S\)-预概形 \(S'\) 及每个由幂零理想 \(\mathscr J\) 定义的子概形 \(S'_0\subset S'\)，从交换方块导出的映射 \[
  F(S')\longrightarrow
  F(S'_0)\times_{G(S'_0)}G(S')
\]

\SGARefTarget{sga3:xi:diagram:1:6:functor-square} \[
\begin{tikzcd}[column sep=huge,row sep=large]
  F(S') \arrow[r] \arrow[d] & G(S') \arrow[d] \\
  F(S'_0) \arrow[r] & G(S'_0)
\end{tikzcd}
\] 是满射（分别为单射、双射）。等价地，对每个 \(S\) 上的函子态射交换方块，其中 \(S'\) 和 \(S'_0\) 如上，\(i:S'_0\to S'\) 为典范浸入，

\SGARefTarget{sga3:xi:diagram:1:6:lifting-square} \[
\begin{tikzcd}[column sep=huge,row sep=large]
  \text{\rm(Q)} & S'_0 \arrow[r,"i"] \arrow[d,"u_0"] &
    S' \arrow[d,"v"'] \\
  & F \arrow[r,"f"] & G ,
\end{tikzcd}
\] 存在至少一个（分别为至多一个、恰好一个）态射

\SGARefTarget{sga3:xi:diagram:1:6:lift} \[
\begin{tikzcd}[column sep=huge,row sep=large]
  & S' \arrow[dl,"u"'] \\
  F &
\end{tikzcd}
\] 使相应的两个三角形交换：\[
  u_0=u\circ i
  \qquad\text{且}\qquad
  v=f\circ u.
\]。

\(S\) 上的函子同态 \(f:F\to G\) 的此性质，在所考虑的函子不可表时仍有意义。若此性质成立，则称同态 \(f:F\to G\) 为\emph{形式平滑}（分别为\emph{形式非分歧}、\emph{形式平展}）的。注意，此性质仅依赖于 \(F\) 和 \(G\) 定义的底函子同态 \[
  (\mathbf{Sch})^\circ\longrightarrow(\mathbf{Ens})
\]（参见 \SGARefLink{sga3:i:proposition:1:4:1}{第 I 篇，1.4.1}），而不依赖于结构态射 \(F\to S\) 和 \(G\to S\)。上述定义有一种在应用中稍方便的等价表述：\(S\) 上的函子态射 \(F\to G\) 形式平滑（分别为形式非分歧、形式平展），若对每个 \(S\) 上的 \(S'\) 及每个态射 \(S'\to G\)，\(S'\) 上的函子 \[
  F'=F\times_G S'
\] 均形式平滑（分别为形式非分歧、形式平展）。
